\exercisegroup

\begin{enumerate}
\def\labelenumi{\arabic{enumi})}
\item
  Let \(\mathrm{F}, \mathrm{G}\) be two vector spaces in duality, such that \(\sigma(\mathrm{F}, \mathrm{G})\) is Hausdorff. Show that if \(\mathcal{T}\) is a Hausdorff topology compatible with the vector space structure of \(\mathrm{F}\) and coarser than \(\sigma(\mathrm{F}, \mathrm{G})\) (but not necessarily locally convex a priori), then \(\mathcal{T} = \sigma(\mathrm{F}, \mathrm{G}_1)\) , where \(\mathrm{G}_1\) is a vector subspace of G, dense in the topology \(\sigma(\mathrm{G},\mathrm{F})\) . (Consider on F the locally convex topology \(T_{1}\) in which a fundamental system of neighbourhoods of 0 is formed by the closed, convex, balanced sets in T which are neighbourhoods of 0 for \(\sigma(\mathrm{F},\mathrm{G})\) .) Deduce that if \(T_{0}\) is a Hausdorff locally convex topology on a vector space E, that is minimal in the set of Hausdorff locally convex topologies on E (II, p.~85, exerc. 13) it is also minimal in the set of topologies (locally convex or not) that are Hausdorff and compatible with the vector space structure of E.
\item
  In \(\mathbf{R}^n\) , let \((\mathrm{C}_i)_{1\leqslant i\leqslant m}\) be a family of \(m\geqslant n + 1\) convex cones with vertex 0; show that if, for any \(n + 1\) of these cones there exists a hyperplane \(\mathbf{H}\) through 0 and such that the cones lie on the same side of \(\mathbf{H}\) , then there exists a hyperplane \(\mathrm{H}_0\) such that all the cones \(\mathrm{C}_i(1\leqslant i\leqslant m)\) lie on the same side of \(\mathrm{H}_0\) (cf.~II, p.~68, exerc. 21, a)).
\item
  In \(\mathbf{R}^n\) let \((\mathrm{D}_i)_{1\leqslant i\leqslant m}\) be a family of \(m\geqslant 2n\) closed half-spaces determined by hyperplanes passing through 0. Show that if, for any \(2n\) of these half spaces, there exists a point \(\neq 0\) in their intersection, then there exists a point \(\neq 0\) in the intersection of all the \(\mathrm{D}_i(1\leqslant i\leqslant m)\) (cf.~II, p.~66, exerc. 10).
\item
  Let S, T be two finite sets in \(R^{n}\) , without common points, such that their union contains at least \(2n + 2\) points. Then there exists a hyperplane separating S from T if, and only if, for every finite set \(F \subset S \cup T\) of \(2n + 2\) points, there exists a hyperplane separating \(F \cap S\) from \(F \cap T\) (use exerc. 3 and the method of II, p.~78, exerc. 15).
\item
  Let E be the vector space of quadratic forms on \(R^{n}\) , which is identified with the vector subspace of symmetric square matrices in the space \(\mathbf{M}_{n}(\mathbf{R})\) of square matrices of order n on R. We endow \(\mathbf{M}_{n}(\mathbf{R})\) with the scalar product \(\operatorname{Tr}(^{t}\mathbf{X},\mathbf{Y})\) , which enables us to identify it with its dual and similarly for E.
\end{enumerate}

\begin{enumerate}
\def\labelenumi{\alph{enumi})}
\item
  Let \(\mathbf{P} \subset \mathbf{E}\) be the set of quadratic forms for which the matrix has all elements \(\geqslant 0\) , and let \(\mathbf{S} \subset \mathbf{E}\) be the set of positive quadratic forms in \(\mathbf{R}^n\) . Show that we have \(\mathbf{P} = \mathbf{P}^\circ\) and \(\mathbf{S} = \mathbf{S}^\circ\) .
\item
  Let B be the set of quadratic forms on \(\mathbf{R}^n\) that can be written in the form \(\sum_{j=1}^{m} x_j'^2\) for same \(m\) , where \(x_{j}^{\prime}\) is a linear form that takes values \(\geqslant 0\) for all \(x = (x_{i})_{1\leqslant i\leqslant n}\) of coordinates \(x_{i}\) all \(\geqslant 0\) ; let \(C\) be the set of quadratic forms that are \(\geqslant 0\) for all the vectors \(x = (x_{i})\) with coordinates \(x_{i}\) all \(\geqslant 0\) . Show that \(B = C^{\circ}\) and \(C = B^{\circ}\) (prove that \(B\) is closed, showing that every element of \(B\) can be written in the form \(\sum_{j=1}^{m} x_{j}^{\prime 2}\) , with \(x_{j}^{\prime}\) positive for all \(x\) with coordinates \(\geqslant 0\) , and \(m \leqslant 2^{n}\) ).
\end{enumerate}

\begin{enumerate}
\def\labelenumi{\arabic{enumi})}
\setcounter{enumi}{5}
\tightlist
\item
  Let F, G be two vector spaces in separating duality, and A a weakly compact convex set in F. Let C be a convex cone with vertex 0, that is weakly closed in G. Suppose that, for all \(y \in \mathbf{C}\) , there exists \(x \in \mathbf{A}\) such that \(\langle x, y \rangle \geqslant 0\) . Show that there exists \(x_0 \in \mathbf{A}\) such that \(\langle x_0, y \rangle \geqslant 0\) for all \(y \in \mathbf{C}\) (apply prop. 4 of II, p.~38, to A and \(\mathbf{C}^\circ\) ).
\end{enumerate}

¶ 7) a) Let F, G be two vector spaces in separating duality and C a weakly closed convex cone in F. Let M be a finite dimensional vector subspace of G. Show that, either there exists \(y_0 \in \mathbf{C}\) such that \(y_0 \in \mathbf{M}^\circ\) and \(y_0 \neq 0\) , or there exists \(z_0 \in \mathbf{M}\) such that \(z_0 \in \mathbf{C}^\circ\) and \(z_0 \neq 0\) (argue by induction on the dimension of M). If C does not contain any line and if the two preceding properties are simultaneously satisfied, show that \(z_0\) cannot be an internal point of \(\mathbf{C}^\circ\) . b) Let the two matrices \((a_{ij}), (b_{ij})\) with real entries in \(n\) rows and \(m\) columns, be such that \(a_{ij} > 0\) for every pair \((i,j)\) . Show that there is a unique value of \(\lambda \in \mathbf{R}\) such that there are two vectors \(x = (x_j) \in \mathbf{R}^m\) , \(y = (y_i) \in \mathbf{R}^n\) satisfying the relations \(x \neq 0, y \neq 0, x_j \geqslant 0, y_i \geqslant 0\) for all \(i,j\) and finally such that

\[
\lambda \sum_ {j = 1} ^ {m} a _ {i j} x _ {j} \geqslant \sum_ {j = 1} ^ {m} b _ {i j} x _ {j} \quad \text { for } \quad 1 \leqslant i \leqslant n\tag{1}
\]

\[
\lambda \sum_ {i = 1} ^ {n} a _ {i j} y _ {i} \leqslant \sum_ {i = 1} ^ {n} b _ {i j} y _ {i} \quad \text { for } \quad 1 \leqslant j \leqslant m\tag{2}.
\]

(Putting \(c_{ij} = \lambda a_{ij} - b_{ij}\) for \(1 \leqslant i \leqslant n\) , \(1 \leqslant j \leqslant m\) , and \(c_{n + i,j} = \delta_{ij}\) (Kronecker's index) for \(1 \leqslant i \leqslant m\) show that the problem reduces to finding a vector \(x \in \mathbf{R}^m\) and a vector \(z = (z_i) \in \mathbf{R}^{n + m}\) which are non null and satisfy the relations

\[
\sum_ {j = 1} ^ {m} c _ {i j} x _ {j} \geqslant 0 \quad \text { for } \quad 1 \leqslant i \leqslant n + m\tag{3}
\]

\[
\sum_ {i = 1} ^ {n + m} c _ {i j} z _ {i} = 0 \quad \text { for } \quad 1 \leqslant j \leqslant m\tag{4}
\]

and \(z_{i} \geqslant 0\) for \(1 \leqslant i \leqslant n + m\) . Remark that, if (3) has a solution for one value \(\lambda_0\) of \(\lambda\) , then it also has a solution for \(\lambda \geqslant \lambda_0\) , and that if (4) has a solution for \(\lambda_0\) , then it also has a solution for \(\lambda \leqslant \lambda_0\) . Finally use \(a)\) .

\begin{enumerate}
\def\labelenumi{\arabic{enumi})}
\setcounter{enumi}{7}
\tightlist
\item
  Let \(\mathbf{T}\) be a compact space and \(\mathbf{L}\) a vector subspace of \(\mathcal{C}(\mathbf{T};\mathbf{R})\) , that is of finite dimension \(r\) ; give to \(\mathbf{L}\) the norm induced by that of \(\mathcal{C}(\mathbf{T};\mathbf{R})\) and to its dual \(\mathbf{L}^*\) the norm \(\| x'\| = \sup_{\| x\| \leqslant 1} \langle x, x' \rangle\) , so that if \(\mathbf{B}\) is the ball \(\| x\| \leqslant 1\) in \(\mathbf{L}\) , then \(\mathbf{B}^\circ\) is the ball \(\| x'\| \leqslant 1\) in \(\mathbf{L}^*\) .
\end{enumerate}

\begin{enumerate}
\def\labelenumi{\alph{enumi})}
\tightlist
\item
  For all \(t \in \mathbf{T}\) , write \(e_t'\) for the linear form \(x \mapsto x(t)\) on \(\mathbf{L}\) . Show that \(\mathbf{B}^\circ\) is the convex envelope of the set of the \(\pm e_t'\) , where \(t\) varies in \(\mathbf{T}\) (obtain a contradiction using prop. 4 of II, p.~38).
\end{enumerate}

Deduce that every linear form \(x' \in \mathbf{L}^*\) such that \(\| x' \| = 1\) we can write \(x' = \sum_{i=1}^{r} \lambda_i e_{t_i}'\) , where

the \(t_i\) are \(r\) points of \(\mathbf{T}\) and the \(\lambda_i\) are real numbers such that \(\sum_{i=1}^{r} |\lambda_i| = 1\) (cf.~II, p.~66, exerc. 9, a)).

\begin{enumerate}
\def\labelenumi{\alph{enumi})}
\setcounter{enumi}{1}
\tightlist
\item
  For every \(y \in \mathcal{C}(\mathbf{T}; \mathbf{R})\) , there exists a unique \(x \in \mathbf{L}\) such that \(\| y - x \| = d(y, \mathbf{L})\) , if and only if for every non null \(z \in \mathbf{L}\) , there exist at most \(r - 1\) distinct points \(t_i \in \mathbf{T}\) such that \(z(t_i) = 0\) (Haar's th.). (To show that this condition is sufficient, observe first that it is equivalent to saying that for \(r\) distinct points \(t_i \in \mathbf{T}\) ( \(1 \leqslant i \leqslant r\) ) the \(e_{t_i}'\) are linearly independent in \(\mathbf{L}^*\) . Now argue by assuming the conclusion is false and obtaining a contradiction. If there exist two distinct points \(x', x''\) of \(\mathbf{L}\) such that \(\| y - x' \| = \| y - x'' \| = d(y, \mathbf{L})\) then there exists \(x_0 \in \mathbf{L}\) and \(z \in \mathbf{L}\) such that, for all sufficiently small real \(\lambda\) , we have \(\| y - (x_0 + \lambda z) \| = d(y, \mathbf{L})\) . Apply the last part of \(a\) ) to the subspace \(\mathbf{L} \oplus \mathbf{R}y\) of \(\mathcal{C}(\mathbf{T}; \mathbf{R})\) and to a suitable linear form on this space which vanishes in \(\mathbf{L}\) . To see that the condition is necessary, note that if it is not true, then there exist \(r\) distinct points \(t_i \in \mathbf{T}\) ( \(1 \leqslant i \leqslant r\) ) such that the \(e_{t_i}'\) are linearly dependent and there exists a function \(z \in \mathbf{L}\) that is not null and vanishes at the points \(t_i\) . If \(\alpha_i (1 \leqslant i \leqslant r)\) are numbers not all zero such that \(\sum_{i=1}^{r} \alpha_i e_{t_i}' = 0\) , consider a function \(w \in \mathcal{C}(\mathbf{T}; \mathbf{R})\) such that \(\| w \| = 1\) , \(w(t_i) = \text{sgn}(\alpha_i)\) for \(1 \leqslant i \leqslant r\) , and the function \(y = w(1 - |\beta z|)\) with \(|\beta|\) sufficiently small and \(\neq 0\) .)
\item
  Suppose that \(T\) is a compact interval in \(R\) and that \(L\) satisfies the condition of Haar's th.; let \((t_i)_{1 \leqslant i \leqslant r+1}\) be a strictly increasing sequence of \(r + 1\) points of \(T\) ; then there exist \(r + 1\) real non zero numbers \(\lambda_i\) such that \(\sum_{i=1}^{r+1} \lambda_i e_{t_i}' = 0\) . Show that then \(\text{sgn}(\lambda_i) \text{sgn}(\lambda_{i+1}) = -1\) for \(1 \leqslant i \leqslant r\) . (Consider separately the case \(r = 1\) and the case \(r > 1\) . In the second case, suppose on the contrary, that for some index \(i \leqslant r - 1\) , the number \(\lambda_i\) is of the same sign as \(\lambda_{i-1}\) or as \(\lambda_{i+1}\) and that \(\lambda_{i-1}\) and \(\lambda_{i+1}\) are of opposite signs. If, for example, \(\lambda_i > 0\) , take \(\alpha_{i-1} > 0\) , \(\alpha_{i+1} > 0\) such that \(\alpha_{i-1} \lambda_{i-1} + \alpha_{i+1} \lambda_{i+1} = 0\) , then \(z \in L\) such that \(z(t_{i-1}) = \alpha_{i-1}\) , \(z(t_{i+1}) = \alpha_{i+1}\) and \(z(t_j) = 0\) for \(j\) distinct from \(i - 1, i\) and \(i + 1\) . Deduce that \(z(t_i) < 0\) and show that this contradicts the given hypotheses.)
\end{enumerate}

\(d)\) With the same hypotheses and notations as those of \(c\) ), suppose that there exists \(y \in \mathcal{C}(\mathrm{T}; \mathbf{R})\) and \(z \in \mathrm{L}\) such that \(y(t_i) - z(t_i) = (-1)^i \alpha_i\) with \(\alpha_i > 0\) for \(1 \leqslant i \leqslant r + 1\) . Show that we then have \(d(y, \mathrm{L}) \geqslant \inf \alpha_i\) . (Use \(a\) ) applied to \(\mathrm{L} \oplus \mathbf{R}y\) , and \(c\) .)

\begin{enumerate}
\def\labelenumi{\alph{enumi})}
\setcounter{enumi}{4}
\tightlist
\item
  With the hypotheses of \(c\) let \(y \in \mathcal{C}(\mathbf{T};\mathbf{R})\) , and \(z\) be the unique point of \(\mathbf{L}\) such that \(\| y - z \| = d(y,\mathbf{L})\) . Show that there exists a strictly increasing sequence \((t_i)_{1\leqslant i\leqslant r + 1}\) in \(\mathbf{T}\) such that
\end{enumerate}

\[
y (t _ {i}) - z (t _ {i}) = (- 1) ^ {i} \varepsilon \| y - z \|
\]

with \(\varepsilon = \pm 1\) . Conversely if \(z\) has this property, then \(z\) is the unique point of \(\mathbf{L}\) such that \(\| y - z \| = d(y, \mathbf{L})\) . (Use \(c\) ) and \(d\) ). Consider the case when \(\mathbf{T}\) is an interval of \(\mathbf{R}\) and when \(\mathbf{L}\) is the set of the restrictions to \(\mathbf{T}\) of polynomials of degree \(< r\) ( \(Tchebycheffs\) th.).

\begin{enumerate}
\def\labelenumi{\arabic{enumi})}
\setcounter{enumi}{8}
\tightlist
\item
  Let \(\mathbf{F}\) and \(\mathbf{G}\) be two vector spaces in separating duality and \(\mathbf{A}\) be a convex subset of \(\mathbf{F}\) which contains 0. For every \(y \in \mathbf{G}\) , write
\end{enumerate}

\[
\mathrm{H} _ {\mathbf {A}} (y) = \sup _ {x \in \mathbf {A}} (- \langle x, y \rangle),
\]

so that \(0 \leqslant \mathrm{H}_{\mathbf{A}}(y) \leqslant +\infty\) ; we call \(\mathrm{H}_{\mathbf{A}}\) the support function of \(\mathbf{A}\) .

\begin{enumerate}
\def\labelenumi{\alph{enumi})}
\item
  Show that \(\mathbf{H}_{\mathrm{A}}\) is the gauge of \(\mathbf{A}^{\circ}\) (II, p.~20).
\item
  If A is weakly compact, then, for all \(y \in \mathbf{G}\) , the hyperplane with the equation \(\langle x, y \rangle = \mathrm{H}_{\mathbf{A}}(y)\) is a support hyperplane of A.
\item
  \(\mathbf{H}_{\mathrm{A}}\) is finite and continuous for the topology \(\sigma (\mathbf{G},\mathbf{F})\) if, and only if, A is finite-dimensional and bounded (in the finite dimensional vector subspace that it generates).
\end{enumerate}

\(d\) ) Let \(\mathbf{A}_i(1\leqslant i\leqslant p)\) be convex sets in F which contain 0 and \(\lambda_{i}\) be real numbers \(\geqslant 0(1\leqslant i\leqslant p)\) ; show that the support function of the convex set \(\mathbf{A} = \sum_{i}\lambda_{i}\mathbf{A}_{i}\) is \(\mathrm{H}_{\mathrm{A}} = \sum_{i}\lambda_{i}\mathrm{H}_{\mathrm{A}_{i}}\) . If \(y\in \mathbf{G}\) is

such that the intersection \(C_{i}\) , of \(A_{i}\) with the hyperplane \(\langle x, y \rangle = H_{A}(y)\) , is non-empty for \(1 \leqslant i \leqslant p\) , show that the intersection of A and the hyperplane with equation \(\langle x, y \rangle = H_{A}(y)\) is the set \(\sum \lambda_{i}C_{i}\) .

\begin{enumerate}
\def\labelenumi{\alph{enumi})}
\setcounter{enumi}{4}
\item
  Suppose that \(\mathbf{A}\) is locally compact and does not contain any line. Then the set which is the union of \(\{0\}\) and the set of \(y \neq 0\) such that \(\mathrm{H}_{\mathrm{A}}(y) = +\infty\) , \(\mathrm{H}_{\mathrm{A}}(-y) \neq +\infty\) is the polar cone of the asymptotic cone \(\mathrm{C}_{\mathrm{A}}\) (consider the case when \(\mathbf{A}\) is the convex envelope of \(\{0\}\) and of one half-line).
\item
  Suppose that F is finite dimensional. Show that, A is parabolic (II, p.~67, exerc. 17) if and only if \(H_{A}\) is a continuous mapping of G in \(\overline{R}\) (if there exists a line parallel to a half-line of \(C_{A}\) , which does not meet A, note that there exists a hyperplane separating this half-line from A).
\end{enumerate}

\begin{enumerate}
\def\labelenumi{\arabic{enumi})}
\setcounter{enumi}{9}
\tightlist
\item
  To each compact convex set A in \(E = R^{n}\) containing 0, we make correspond its support function \(H_{A}\) by the duality between E and \(E^{*} : H_{A}\) belonging to the space \(\mathcal{C}(E^{*}; \mathbf{R})\) of continuous real valued functions in \(E^{*}\) . We ascribe to the space \(\mathcal{C}(E^{*}; \mathbf{R})\) the uniform structure of compact convergence and, to the set \(\Re_{0}^{\prime}(E)\) of the compact convex sets in E containing 0, the uniform structure defined in the exerc. 39 of II, p.~71. Show that \(A \mapsto H_{A}\) is an isomorphism of \(\Re_{0}^{\prime}(E)\) on a uniform subspace of \(\mathcal{C}(E^{*}; \mathbf{R})\) .
\end{enumerate}

Deduce that the mapping \(\mathbf{A} \mapsto \mathbf{A}^{\circ}\) of the set \(\Re_0(\mathbf{E})\) of compact convex sets in \(\mathbf{E}\) which contain 0 as an interior point, on the set \(\Re_0(\mathbf{E}^*)\) , is an isomorphism for the uniform structures of these two spaces (cf.~II, p.~71, exerc. 39).

\begin{enumerate}
\def\labelenumi{\arabic{enumi})}
\setcounter{enumi}{10}
\tightlist
\item
  Let \(\mathrm{F}, \mathrm{G}\) be two vector spaces in separating duality. An ultrafilter \(\mathfrak{U}\) on \(\mathrm{F}\) converges weakly to a point \(x_0\) if, and only if, \(x_0\) belongs to the intersection of all the weakly closed convex sets which belong to \(\mathfrak{U}\) (note that if \(x_0\) is a point of this intersection that is not a cluster point of \(\mathfrak{U}\) , then there exists a closed half-space belonging to \(\mathfrak{U}\) and not containing \(x_0\) ).
\end{enumerate}

Deduce from this result that, for a sequence of points \((x_{n})\) of F to be weakly convergent to a point a, it is necessary and sufficient that a belongs to all the weakly closed convex envelopes of the sets formed by an infinity of the terms of the sequence (use prop. 7 of GT, I, § 6.4).

\begin{enumerate}
\def\labelenumi{\arabic{enumi})}
\setcounter{enumi}{11}
\tightlist
\item
  \begin{enumerate}
  \def\labelenumii{\alph{enumii})}
  \tightlist
  \item
    Let E be a vector space and \((\mathrm{E}_{\alpha})_{\alpha\in\mathrm{A}}\) be an increasing directed family of subspaces of E, whose union is E; each \(E_{\alpha}\) is supposed to carry a locally convex topology \(T_{\alpha}\) such that for \(\alpha\leqslant\beta\) the canonical injection \(E_{\alpha}\to E_{\beta}\) is continuous. Let T be the topology on E, which is the inductive limit of the \(T_{\alpha}\) (II, p.~29, Example II); show that the dual \(E'\) of E (for T) with the topology \(\sigma(\mathrm{E}',\mathrm{E})\) can be canonically identified with the projective limit of the duals \(E'_{\alpha}\) , with topology \(\sigma(\mathrm{E}'_{\alpha},\mathrm{E}_{\alpha})\) .
  \end{enumerate}
\end{enumerate}

\begin{enumerate}
\def\labelenumi{\alph{enumi})}
\setcounter{enumi}{1}
\tightlist
\item
  Let \((\mathbf{X}_{\alpha},\phi_{\alpha \beta})\) be a projective system of non-empty sets corresponding to a directed set of indices A, such that \(\phi_{\alpha \beta}\) are surjective and that \(\varprojlim \mathbf{X}_{\alpha} = \emptyset\) \((\mathbf{S},\mathbf{III},\S 7,\text{exerc. }4)\) . Put \(\mathbf{F}_{\alpha} = \mathbf{R}^{(\mathbf{X}_{\alpha})}\) and denote by \(f_{\alpha \beta}:\mathrm{F}_{\beta}\to \mathrm{F}_{\alpha}\) for \(\alpha \leqslant \beta\) the linear mapping deduced canonically from \(\phi_{\alpha \beta}\) (A, II, § 1.11, cor. 1). If, we give to each \(\mathbf{F}_{\alpha}\) the topology which is the direct sum topology of its factors, the dual \(\mathrm{E}_{\alpha} = \mathrm{F}_{\alpha}'\) of \(\mathrm{F}_{\alpha}\) , with the weak topology \(\sigma(\mathrm{E}_{\alpha}, \mathrm{F}_{\alpha})\) , can be identified with the product space \(\mathbf{R}^{\mathbf{X}_{\alpha}}\) , and \({}^{t}f_{\alpha\beta}\) is an isomorphism of \(\mathrm{E}_{\alpha}\) on a closed subspace of \(\mathrm{E}_{\beta}\) , having a topological complement in \(\mathrm{E}_{\beta}\) . Show that on \(\mathrm{E} = \varinjlim \mathrm{E}_{\alpha}\) (for the \({}^{t}f_{\alpha\beta}\) ) the topology which is the inductive limit of those of \(\mathrm{E}_{\alpha}\) is the coarsest topology (therefore non-Hausdorff) (using \(a\) ) and noting that \(\varprojlim \mathrm{F}_{\alpha} = \{0\}\) ).
\end{enumerate}

\begin{enumerate}
\def\labelenumi{\arabic{enumi})}
\setcounter{enumi}{12}
\tightlist
\item
  Let E be a vector space. We say that a Hausdorff locally convex topology T on E is minimal (and that E, with T, is a space of minimal type) if there exists no Hausdorff locally convex topology on E, that is strictly coarser than T (cf.~II, p.~81, exerc. 1).
\end{enumerate}

\begin{enumerate}
\def\labelenumi{\alph{enumi})}
\item
  Let \(\mathcal{T}\) be a minimal topology on E, and let \(E'\) be the dual of E (when E has the topology \(\mathcal{T}\) ); show that \(\mathcal{T} = \sigma(E, E')\) and \(E = E'^*\) (note that there cannot be an everywhere dense hyperplane in \(E'\) for the topology \(\sigma(E', E)\) using the cor. 3 of II, p.~43). Deduce that spaces of minimal type are products of lines.
\item
  Show that in a Hausdorff locally convex space F, every subspace E of minimal type has a topological complement, and in particular is closed (use a) and the Hahn-Banach th. for extending the identity mapping of E in itself to a mapping of F in E).
\item
  Let \(u\) be a continuous linear mapping of a space \(\mathbf{E}\) of minimal type in a Hausdorff locally convex space \(\mathbf{F}\) . Show that \(u(\mathbf{E})\) is closed in \(\mathbf{F}\) and that \(u\) is a strict morphism of \(\mathbf{E}\) in \(\mathbf{F}\) (use \(b\) ) and the definition of a space of minimal type).
\item
  Let F be a Hausdorff locally convex space and M be a closed vector subspace of F. Show that, if there exists a complement N of M in F that is a subspace of minimal type, then N is a topological complement of M in F (use c)).
\item
  Let M be a subspace of minimal type in a Hausdorff locally convex space F; show that, for every closed vector subspace N of F the sum \(M + N\) is closed in F (consider the quotient space F/N and use c)). If further N is of minimal type, then \(M + N\) is of minimal type.
\end{enumerate}

\begin{enumerate}
\def\labelenumi{\arabic{enumi})}
\setcounter{enumi}{13}
\tightlist
\item
  Let \(\mathbf{E}\) be a Hausdorff locally convex space and \(\mathbf{F}\) a locally convex space of minimal type (exerc. 13).
\end{enumerate}

\begin{enumerate}
\def\labelenumi{\alph{enumi})}
\item
  Show that if \(\mathbf{M}\) is a closed vector subspace of the product space \(\mathrm{E} \times \mathrm{F}\) , its projection on \(\mathrm{E}\) is closed in \(\mathrm{E}\) (use exerc. 13, e)).
\item
  Let \(u\) be a linear mapping of \(\mathbf{E}\) in \(\mathbf{F}\) . Show that if the graph of \(u\) is closed in \(\mathbf{E} \times \mathbf{F}\) , then \(u\) is continuous (use \(a\) ).
\item
  Suppose that, in E, every closed vector subspace has a topological complement (cf.~V, p.~13). Show that, in \(E \times F\) , every closed vector subspace M has a topological complement. (If \(N_{1}\) is the projection of M on E and \(N_{2}\) a topological complement of \(N_{1}\) in E, if \(P_{1} = M \cap F\) , and \(P_{2}\) is a topological complement of \(P_{1}\) in F, show that \(N_{2} + P_{2}\) is a topological complement of M in \(E \times F\) , using b).)
\end{enumerate}

* 15) Let E, F be two Hausdorff locally convex space. We say that a continuous linear mapping \(u: \mathrm{E} \to \mathrm{F}\) is linearly proper if, for every Hausdorff locally convex space G and every closed vector subspace V of \(\mathrm{E} \times \mathrm{G}\) the image of V by \(u \times 1_{\mathrm{G}}: \mathrm{E} \times \mathrm{G} \to \mathrm{F} \times \mathrm{G}\) is closed. Show that this condition is equivalent to the following: \(u^{-1}(0)\) is a subspace of minimal type of E and for every closed vector subspace W of E, the set \(u(\mathbf{W})\) is closed in F. (To show that the first condition implies the second, consider the mapping \(v: \mathrm{E} \to \{0\}\) and, giving E the topology \(\sigma(\mathrm{E}, \mathrm{E}')\) , so that E is immersed in \(\mathrm{E}'*\) with \(\sigma(\mathrm{E}'*, \mathrm{E}')\) , take the image under the projection \(v \times 1_{\mathrm{E}'*}: \mathrm{E} \times \mathrm{E}'* \to \mathrm{E}'*\) of the closure in \(\mathrm{E} \times \mathrm{E}'*\) of the diagonal \(\Delta\) of \(\mathrm{E} \times \mathrm{E}\) . To show that the second condition implies the first, show that it implies that, for the topologies \(\sigma(\mathrm{E}, \mathrm{E}')\) and \(\sigma(\mathrm{F}, \mathrm{F}')\) , the mapping \(u\) is a strict morphism and use exerc. 13, e).) *

\begin{enumerate}
\def\labelenumi{\arabic{enumi})}
\setcounter{enumi}{15}
\tightlist
\item
  Let \(\mathbf{F}\) be a product of lines and \(\mathbf{C}\) a closed convex set in \(\mathbf{F}\) .
\end{enumerate}

\begin{enumerate}
\def\labelenumi{\alph{enumi})}
\item
  Show that there exists \(x_0 \in \mathbf{F}\) , two sets I and J and a topological isomorphism \(u\) of \(\mathbf{F}\) on \(\mathbf{R}^{\mathbf{I}} \times \mathbf{R}^{\mathbf{J}}\) such that \(u(x_0 + \mathbf{C})\) is of the form \(\mathbf{R}^{\mathbf{I}} \times \mathbf{A}\) , where \(\mathbf{A}\) is a closed convex set of \(\mathbf{R}_+^{\mathbf{J}}\) . (Note that we have \(\mathbf{F} = \mathbf{G}^*\) where \(\mathbf{F}\) has the topology \(\sigma(\mathbf{G}^*, \mathbf{G})\) ; consider the polar \(C^\circ\) of \(C\) in \(G\) , the vector subspace of \(G\) generated by \(C^\circ\) and a complement of this subspace.)
\item
  If C does not contain any affine line, the mapping \((x, y) \mapsto x + y\) of \(C \times C\) in \(F\) is proper.
\item
  Suppose that C is a cone with vertex 0 and that the uniform structure induced on C by that of F is metrisable. Then if the sets I and J, the point \(x_0\) and the mapping \(u\) satisfy the conditions of \(a\) ) then I is enumerable and there exists an enumerable subset H of J such that the restriction of the canonical projection \(p: \mathbf{R}^{\mathrm{I}} \times \mathbf{R}^{\mathrm{J}} \to \mathbf{R}^{\mathrm{I}} \times \mathbf{R}^{\mathrm{H}}\) to \(u(x_0 + \mathbf{C})\) is an isomorphism of the uniform subspace \(u(x_0 + \mathbf{C})\) of \(\mathbf{R}^{\mathrm{I}} \times \mathbf{R}^{\mathrm{J}}\) on the uniform subspace \(p(u(x_0 + \mathbf{C}))\) of \(\mathbf{R}^{\mathrm{I}} \times \mathbf{R}^{\mathrm{H}}\) .
\end{enumerate}

\begin{enumerate}
\def\labelenumi{\arabic{enumi})}
\setcounter{enumi}{16}
\tightlist
\item
  Let \(\mathbf{E}\) be an infinite dimensional vector space.
\end{enumerate}

\begin{enumerate}
\def\labelenumi{\alph{enumi})}
\tightlist
\item
  Show that there exist hyperplanes in \(\mathrm{E}^*\) that are everywhere dense for the topology \(\sigma (\mathrm{E}^{*},\mathrm{E})\) b) If \(\mathbf{H}'\) is such a hyperplane, show that, in E, the only linear subvarieties \(\neq \mathrm{E}\) that are everywhere dense for the topology \(\sigma (\mathrm{E},\mathrm{H}')\) are the hyperplanes.
\end{enumerate}

  ¶ 18) a) In a normed space E, let A be a closed convex set \(\neq\) E; show that the function \(x\mapsto d(x,\complement\mathbf{A})\) is concave in A (use the fact that A is the intersection of closed half-spaces). b) Define inductively a sequence of closed convex sets \(\mathbf{A}_{n}\subset\mathbf{R}^{n}\) in the following manner; \(\mathbf{A}_{1}=\mathbf{R}_{+}\) ; if \(\mathbf{R}^{n+1}\) is identified with \(\mathbf{R}^{n}\times\mathbf{R}\) , then \(\mathbf{A}_{n+1}\) is the set of pairs \((x,\xi)\) such that \(x\in\overset{\circ}{\mathbf{A}}_{n}\) and that

\[
\xi \geqslant (d (x, \complement A _ {n})) ^ {- 1} + \| x \| ^ {2},
\]

where \(\|x\|\) is the Euclidean norm. Show that \(A_{n+1}\) does not have any support hyperplane of the form \(H \times R\) , where H is a hyperplane of \(R^{n}\) and that its asymptotic cone is \(\{0\} \times R_{+}\) . c) If \(p_{nm}\) is the canonical projection \(R^{m} \to R^{n}\) ( \(R^{m}\) being identified with \(R^{n} \times R^{m-n}\) ) for \(m \geqslant n\) , show that when \(R^{N}\) is identified with the projective limit of the projective system ( \(R^{n}, p_{nm}\) ) the \(A_{n}\) form a projective system of sets and that \(A = \varprojlim A_{n}\) is a closed convex set not relatively compact in \(R^{N}\) , having no closed hyperplane of support and such that \(C_{A} = \{0\}\) .

\begin{enumerate}
\def\labelenumi{\arabic{enumi})}
\setcounter{enumi}{18}
\tightlist
\item
  \begin{enumerate}
  \def\labelenumii{\alph{enumii})}
  \tightlist
  \item
    Let A be a closed convex set in E, a product of lines, that is non-compact and such that \(C_{A} = \{0\}\) (exerc. 18). Show that if \(B = \overline{A - A}\) and if M is the convex closed envelope of \(A \cup (-A)\) , then B and M contain lines (use exerc. 16, b) of II, p.~85).
  \end{enumerate}
\end{enumerate}

\begin{enumerate}
\def\labelenumi{\alph{enumi})}
\setcounter{enumi}{1}
\tightlist
\item
  Let \(\mathrm{A}_1, \mathrm{A}_2\) be two closed convex sets in E such that \(\mathrm{A}_1 + \mathrm{A}_2\) is closed and none of \(\mathrm{A}_1\) , \(\mathrm{A}_2, \mathrm{A}_1 + \mathrm{A}_2\) contain affine lines. Show that \(\mathrm{C}_{\mathrm{A}_1 + \mathrm{A}_2} = \mathrm{C}_{\mathrm{A}_1} + \mathrm{C}_{\mathrm{A}_2}\) (use exerc. 16, b) of II, p.~85). c) Let A be a closed convex set in E that does not contain any affine line and \(\mathbf{M}_1, \dots, \mathbf{M}_n\) closed convex sets contained in A. If B is the convex envelope of \(\bigcup_{i} \mathbf{M}_i\) , show that \(\overline{\mathbf{B}} = \mathbf{B} + \sum_{i} \mathbf{C}_{\mathbf{M}_i}\)
\end{enumerate}

and \(C_B = \sum_i C_{M_i}\) (same method).

¶ 20) Let \(\mathrm{F} = \mathbf{R}^{(\mathrm{A})}, \mathrm{G} = \mathbf{R}^{\mathrm{A}}\) , where \(\mathrm{A}\) is any infinite set; suppose that \(\mathrm{F}\) and \(\mathrm{G}\) are put in separating duality by the bilinear form \(\langle x, y \rangle = \sum_{\alpha \in \mathbf{A}} x(\alpha) y(\alpha)\) .

\begin{enumerate}
\def\labelenumi{\alph{enumi})}
\item
  Let N be an additive subgroup of G; we denote by \(N^{*}\) the subgroup of the \(x \in F\) such that \(\langle x, y \rangle\) is an integer for all \(y \in N\) and by \(N^{**}\) the subgroup of the \(z \in G\) such that \(\langle x, z \rangle\) is an integer for all \(x \in N^{*}\) . If \(\overline{N}\) is the closure of N for the topology \(\sigma(G, F)\) , show that \(N^{*}\) is closed in F for \(\sigma(F, G)\) and that \(N^{**} = \overline{N}\) (to establish this last point, use GT, VII, § 1.3, prop. 6, projecting N on the finite dimensional coordinate varieties of G).
\item
  Suppose that \(\mathbf{A} = \mathbf{N}\) . Let \(\mathbf{M}\) be a closed subgroup of \(\mathbf{F}\) for \(\sigma(\mathbf{F},\mathbf{G})\) ; show that if \(\mathbf{V}\) is the largest vector subspace contained in \(\mathbf{M}\) , then \(\mathbf{M}\) is the topological direct sum of \(\mathbf{V}\) and of a closed subgroup \(\mathbf{P}\) that is a free \(\mathbf{Z}\) -module having an enumerable base. (Consider \(\mathbf{F}\) as the union of an increasing sequence \((\mathrm{F}_n)\) of finite dimensional vector subspace and apply GT, VII, § 1.2, th. 2 and § 1, exerc. 7.) \(\mathbf{P}\) is discrete (for the topology induced by \(\sigma(\mathbf{F},\mathbf{G})\) ), if, and only if, \(\mathbf{P}\) is of finite rank.
\item
  Deduce from a) and b) that when A = N, every closed subgroup of G (when G carries the product topology \(\sigma(\mathbf{G}, \mathbf{F})\) ) can be transformed, by an automorphism of the topological group G, in a product \(R^{I} \times Z^{J}\) , where I and J are two sets of N without common elements. d) In the space \(E = R^{N}\) , carrying the topology \(\sigma(E, E^{*})\) , show that the subgroup \(Z^{N}\) is closed and does not contain any line, even though it is not a free Z-module (A, VII, p.~59, exerc. 8); the results of b) do not therefore extend when A is not enumerable.
\end{enumerate}

